% Original top-view schematic based on IO and Disk.pdf, slides 10.15--10.16.
% An annular sector is highlighted; the head lies on the same track at angle 0.
\begin{tikzpicture}[font=\footnotesize,>=Latex]
  \fill[gray!7] (0,0) circle (1.65);
  \draw[gray] (0,0) circle (1.65);
  \draw[gray!65] (0,0) circle (0.7);
  \draw[noteBlue] (0,0) circle (1.25);
  \filldraw[fill=noteBlue!25,draw=noteBlue]
    (35:1.1) arc (35:65:1.1) -- (65:1.4)
    arc (65:35:1.4) -- cycle;
  \fill[gray!60] (0,0) circle (0.14);
  \draw[very thick] (1.25,0) -- (2.65,0);
  \fill[noteBlue] (1.18,-0.07) rectangle (1.32,0.07);
  \node[anchor=west] at (2.75,0) {磁头及磁臂};
  \node[anchor=west] (sector) at (2.3,1.3) {目标扇区};
  \draw[->] (sector.west) -- (50:1.26);
  \node[anchor=east] (track) at (-2,0.65) {磁道};
  \draw[->] (track.east) -- (150:1.25);
  \draw[->,noteBlue,thick] (85:1.85) arc (85:145:1.85);
  \node[anchor=east] at (-1.1,1.95) {盘片旋转};
  \draw[<->,thick] (0.5,-0.4) -- (2.1,-0.4);
  \node[anchor=west] at (2.3,-0.65) {寻道：沿半径移动};
\end{tikzpicture}
